\thispagestyle{empty}\vspace*{\fill}
\begin{minipage}{128mm}\fontsize{9}{13}\selectfont\color{otMuted}
{\otSemi\color{otInk}Mathématiques · Quatrième · France}\par
{\otSemi\color{otGreenDark}Édition 2027-2028 · Nouveau programme}\par\medskip
Conception et mise en page : collection OptiTAB.\par
270 exercices de chapitre, 18 défis et deux bilans annuels originaux, corrigés.\par
Charte graphique de la collection OptiTAB, du collège au lycée.\par\bigskip
Édition du 23 octobre 2026 · Année scolaire 2027-2028.\par\bigskip
Cette édition suit le nouveau programme de mathématiques du cycle 4
(arrêté du 18 février 2026, Bulletin officiel n\textsuperscript{o} 10 du 5 mars 2026),
applicable en Quatrième à partir de la rentrée 2027. Elle remplace, pour
cette rentrée, l’édition 2026-2027, conforme au programme de 2020.\par\bigskip
L’ordre, les exercices et les barèmes sont des choix éditoriaux. Les défis sont
facultatifs. Les deux bilans sont des évaluations d’entraînement originales,
sans statut d’examen officiel.\par\bigskip
Les données des situations sont construites pour apprendre ; ce ne sont pas
des relevés économiques ou scientifiques. Sauf mention contraire, les
schémas de principe ne sont pas à mesurer. Les coordonnées et graduations
explicites font foi dans les graphiques.\par\bigskip
Conserver les résultats exacts, sauf demande d’arrondi. Les réponses ajoutées
en vert dans les figures corrigées ne sont pas des données de l’énoncé.\par\bigskip
© OptiTAB, 2026.\par\otSite
\end{minipage}\clearpage
